\section{Governing equations}%
\label{sec:governing_equation}
In this work we consider two physical models to describe the Rayleigh--Taylor instability that develops between two fluids of different densities. The first model is the Boussinesq approximation that provides the simplest model for examining the effect of buoyant stratification on the flow. Density variations relative to a reference state are assumed to be small and affect only the gravitational buoyancy term, while the velocity field remains incompressible~\citep{boffetta-mazzino-2017}. The other model that we employ is the incompressible variable-density (IVD) formulation of \citet{sandoval-1995}. It applies in the low-Mach-number regime, where acoustic waves are suppressed and the energy equation decouples from the problem~\citep{livescu-ristorcelli-2007}. Each fluid remains individually incompressible, although the mixture velocity is not divergence-free~\citep{livescu-2013}.

For both models we consider miscible fluids with reference densities \(\rho_1\) and \(\rho_2\)
and dynamic viscosities \(\mu_1\) and \(\mu_2\). We nondimensionalize the governing equations using the reference density
\(\rho_r\), length \(L_r\), velocity \(U_r\), dynamic viscosity \(\mu_r\), gravitational
acceleration \(g\), and mass diffusivity \(\mathcal{D}\). The resulting Reynolds, Froude, and Schmidt numbers
are
\begin{equation}
    \label{eq:Re_Fr_definitions}
    \Rey = \frac{\rho_r L_r U_r}{\mu_r}, \qquad \Frd = \frac{U_r}{\sqrt{g L_r}}, \qquad \Sch = \frac{\mu_r/\rho_r}{\mathcal{D}}.
\end{equation}
The present problem has no externally imposed length or velocity scale, so the reference length and velocity are set by the classical viscous-gravitational scales
of~\citet{chandrasekhar-1961}. With this choice, the Reynolds and Froude numbers reduce to unity, and the reference scales become
\begin{equation}
    \label{eq:chandrasekhar_scaling}
    U_r = (\nu_r g)^{1/3}, \qquad T_r = \left(\frac{\nu_r}{g^{2}}\right)^{1/3}, \qquad L_r = \left(\frac{\nu_r^{2}}{g}\right)^{1/3},
\end{equation}
where \(\nu_r=\mu_r/\rho_r\) is the reference kinematic viscosity. The reference density
and dynamic viscosity are taken as the arithmetic means of the two fluids,
\begin{equation}
    \label{eq:reference_quantities}
    \rho_r = \frac{\rho_1 + \rho_2}{2}, \qquad \mu_r = \frac{\mu_1 + \mu_2}{2},
\end{equation}
and the density and viscosity contrasts are measured by the density Atwood number \(\Atw\) and a corresponding viscosity Atwood number \(A_{\mu}\),
\begin{equation}
    \label{eq:atwood_numbers}
    \Atw = \frac{\rho_2 - \rho_1}{\rho_2 + \rho_1}, \qquad A_{\mu} = \frac{\mu_2 - \mu_1}{\mu_2 + \mu_1}.
\end{equation}
We consider two models of the Rayleigh--Taylor instability. The first is the Boussinesq formulation presented in \cref{sec:bq_equation}. It is valid for weak stratification and provides the simplest setting in which to develop the statistical approach, interpret its predictions, and validate it against three-dimensional direct numerical simulations. The second is the full variable-density formulation presented in \cref{sec:vd_equation}, which retains non-Boussinesq effects throughout the transient instability.

\subsection{Boussinesq approximation}
\label{sec:bq_equation}
In this model, the density is taken to vary linearly with the local mixture concentration, and the governing equations are
\begin{subequations}
    \label{eq:bq_nl_continuous}
    \begin{align}
        \label{eq:bq_nl_continuous_momentum}
        \diff{\boldsymbol{u}}{t} + (\boldsymbol{u}\bcdot\bnabla)\boldsymbol{u} &= -\bnabla p + \nabla^{2}\boldsymbol{u} - c\,\boldsymbol{e}_{z}, \\
        \label{eq:bq_nl_continuous_scalar}
        \diff{c}{t} + \boldsymbol{u}\bcdot\bnabla c &= \Sch^{-1}\nabla^{2}c, \\
        \label{eq:bq_nl_continuous_continuity}
        \bnabla\bcdot\boldsymbol{u} &= 0,
    \end{align}
\end{subequations}
where \(\boldsymbol{u}(\boldsymbol{x},t)\) is the velocity field, \(p(\boldsymbol{x},t)\) is the pressure field, and \(c(\boldsymbol{x},t)\) is the scalar concentration field. The viscosity is taken constant throughout the mixture by setting \(A_{\mu}=0\); the effect of viscosity stratification is deferred to the variable-density model of \cref{sec:vd_equation}, which recovers the Boussinesq approximation in the limit of weak stratification.

Introducing the state vector \(\boldsymbol{Q}(\boldsymbol{x},t) = (u, v, w, p, c)^\top\), the nonlinear system can be written compactly as
\begin{equation}
    \label{eq:compact_nl_continuous}
    \mathscr{M}\,\diff{\boldsymbol{Q}}{t} = \mathcal{N}\!\left( \boldsymbol{Q} \right),
\end{equation}
where \(\mathscr{M}\) is the identity matrix with the entry corresponding to pressure set to zero, since the pressure is not dynamically evolved but instead enforces the incompressibility constraint.
We introduce the perturbation expansion
\begin{equation}
    \label{eq:perturbation_expansion}
    \boldsymbol{Q}(\boldsymbol{x},t) = \boldsymbol{Q}_0(\boldsymbol{x},t) + \epsilon\,\boldsymbol{Q}'(\boldsymbol{x},t) + O(\epsilon^2),
    \qquad 0 < \epsilon \ll 1,
\end{equation}
where \(\boldsymbol{Q}_0\) is a solution of the nonlinear equations, hereafter referred to as the base state, and \(\boldsymbol{Q}'\) is the perturbation about this base state. Substituting this expansion into the compact form and collecting terms of equal order in \(\epsilon\), we obtain the governing equations
\begin{subequations}
    \label{eq:compact_lin_continuous}
    \begin{align}
        \label{eq:compact_lin_continuous_0th}
        \mathscr{M}\,\diff{\boldsymbol{Q}_0}{t} &= \mathcal{N}\!\left( \boldsymbol{Q}_0 \right), \\
        \label{eq:compact_lin_continuous_1th}
        \mathscr{M}\,\diff{\boldsymbol{Q}'}{t} &= \mathcal{L}(\boldsymbol{Q}_0)\,\boldsymbol{Q}',
    \end{align}
\end{subequations}
where \(\mathcal{L}\) is the Jacobian of the nonlinear Boussinesq operator. Since no intrinsic length or time scale is imposed on the system, the base-state evolution is assumed to depend only on the vertical coordinate and time~\citep{prathama-pantano-2021}, leading to \(\boldsymbol{u}_0=0\) and
\begin{equation}
    \label{eq:bq_base_state}
    c_0(z,t) = \frac{1}{2}\left[1-\erf\!\left(\frac{z}{\delta(t)}\right)\right], \qquad \delta(t) = \sqrt{4t/\Sch}.
\end{equation}
The Boussinesq linearized equations can be further simplified by eliminating the pressure term. This can be done easely by taking the difference between the vertical derivative of the poisson equation and the laplacian of the vertical momentum to get rid of the pressure term and exploting the zero divergence constraint. The resulting system for the reduced perturbation state \(\boldsymbol{q}=(w,c)^{\top}\) is
\begin{equation}
    \label{eq:bq_lin_continuous}
    \begin{pmatrix}
        \nabla^{2} & 0 \\
        0 & 1
    \end{pmatrix}
    \begin{pmatrix}
        \partial_t\, w \\[4pt]
        \partial_t\, c
    \end{pmatrix}
    =
    \begin{pmatrix}
        \nabla^{4} & \nabla_{xy}^{2} \\[4pt]
        -\partial_{z}c_{0} & \Sch^{-1}\nabla^{2}
    \end{pmatrix}
    \begin{pmatrix}
        w \\[4pt]
        c
    \end{pmatrix}.
\end{equation}
The boundary conditions enforce vanishing vertical velocity and its vertical gradient at the
domain boundaries, and require the concentration perturbation to vanish in the far field,
\begin{equation}
    \label{eq:bq_boundary_conditions}
    w(\pm L_z) = 0, \qquad \left.\frac{\partial w}{\partial z}\right|_{\pm L_z} = 0, \qquad
    c(\pm L_z) = 0.
\end{equation}
Assuming homogeneity in the horizontal directions \((x,y)\), the perturbation fields are decomposed into Fourier modes. Fourier-transforming the reduced system then yields, for the
state vector \(\hat{\boldsymbol{q}}=(\hat{w},\hat{c})^{\top}\),
\begin{equation}
    \label{eq:lti_dz_k_boussinesq}
    \begin{pmatrix}
        \partial_{zz}-k^{2} & 0 \\
        0 & 1
    \end{pmatrix}
    \begin{pmatrix}
        \partial_t\, \hat{w} \\[4pt]
        \partial_t\, \hat{c}
    \end{pmatrix}
    =
    \begin{pmatrix}
        \left(\partial_{zz}-k^{2}\right)^{2} & -k^{2} \\[4pt]
        -\partial_{z}c_{0} & \Sch^{-1}\left(\partial_{zz}-k^{2}\right)
    \end{pmatrix}
    \begin{pmatrix}
        \hat{w} \\[4pt]
        \hat{c}
    \end{pmatrix},
\end{equation}
where \(\hat{w}(k,z,t)\) and \(\hat{c}(k,z,t)\) are the Fourier-transformed velocity and
concentration perturbations. After discretizing the inhomogeneous vertical direction, the state associated with a given wavenumber is denoted by \(\hat{\mathbf{q}}_{k}(t)\in\mathbb{C}^{2N_z}\), where \(N_z\) is the number of degrees of freedom. The continuous differential operators then become the time-dependent matrices \(\mathsfbi{A}_{k}(t)\) and \(\mathsfbi{B}_{k}(t)\), and the linearized dynamics take the form
\begin{equation}
    \label{eq:semidiscrete_system}
    \mathsfbi{B}_{k}(t)\,\difft{\hat{\mathbf{q}}_{k}}{t}
    =
    \mathsfbi{A}_{k}(t)\,\hat{\mathbf{q}}_{k},
    \qquad
    \hat{\mathbf{q}}_{k}(0)=\hat{\mathbf{q}}_{0,k}.
\end{equation}
This semi-discrete initial-value problem defines the discrete propagator
\(\boldsymbol{\Phi}_{k}(t,0)\) through
\begin{equation}
    \label{eq:discrete_propagator}
    \hat{\mathbf{q}}_{k}(t)
    =
    \boldsymbol{\Phi}_{k}(t,0)\,\hat{\mathbf{q}}_{0,k},
    \qquad
    \boldsymbol{\Phi}_{k}(0,0)=\mathsfbi{I}.
\end{equation}
The two-argument form \(\boldsymbol{\Phi}_{k}(t,0)\) is cumbersome for sustained use. We therefore adopt the condensed notation
\begin{equation}
    \label{eq:propagator_shorthand}
    \boldsymbol{\Phi}_{t,k} \equiv \boldsymbol{\Phi}_{k}(t,0),
\end{equation}
so that the discrete propagator takes the concise form \(\hat{\mathbf{q}}_{k}(t)=\boldsymbol{\Phi}_{t,k}\,\hat{\mathbf{q}}_{0,k}\). This shorthand is used throughout the remainder of this work.

\subsection{Variable-density formulation}
\label{sec:vd_equation}
For the IVD, with the nondimensionalization introduced above, the governing equations are
\begin{subequations}
    \label{eq:vd_nl_continuous}
    \begin{align}
        \label{eq:vd_nl_continuous_continuity}
        \partial_t \rho + \bnabla\bcdot(\rho \boldsymbol{u}) &= 0, \\
        \label{eq:vd_nl_continuous_momentum}
        \rho\,\partial_t \boldsymbol{u} + \rho\,(\boldsymbol{u}\bcdot\bnabla)\boldsymbol{u} &= -\,\bnabla p + \bnabla\bcdot\boldsymbol{\tau} - \rho\,\boldsymbol{e}_z, \\
        \label{eq:vd_nl_continuous_divergence}
        \bnabla \bcdot \boldsymbol{u} &= -\Sch^{-1}\,\nabla^2 \log\rho,
    \end{align}
\end{subequations}
where \(\rho(\boldsymbol{x},t)\) is the mixture density, \(\boldsymbol{u}(\boldsymbol{x},t)\) is the
velocity field, \(p(\boldsymbol{x},t)\) is the pressure, and \(\boldsymbol{e}_z\) points
opposite to gravity. The viscous stress tensor \(\boldsymbol{\tau}\) is given by
\begin{equation}
    \label{eq:vd_stress_tensor}
    \boldsymbol{\tau} =  2\,\mu\,\mathsfbi{S} - \frac{2}{3}\mu\,(\bnabla\bcdot\boldsymbol{u})\mathsfbi{I}, \qquad 2\,\mathsfbi{S} = \bnabla\boldsymbol{u} + (\bnabla\boldsymbol{u})^{\!\top},
\end{equation}
where \(\mu\) is the dynamic viscosity, \(\mathsfbi{S}\) is the strain-rate tensor, and
\(\mathsfbi{I}\) is the identity tensor.

Introducing the state vector \(\boldsymbol{Q}(\boldsymbol{x},t) = (u, v, w, p, \rho)^\top\), the nonlinear system can be written compactly as
\begin{equation}
    \label{eq:vd_compact_nl_continuous}
    \mathscr{M}\,\diff{\boldsymbol{Q}}{t} = \mathcal{N}\!\left( \boldsymbol{Q} \right),
\end{equation}
where \(\mathscr{M}\) is again the identity matrix with the entry corresponding to \(p\) set to zero, since the pressure is not dynamically evolved but instead enforces the divergence constraint.
We introduce the perturbation expansion
\begin{equation}
    \label{eq:vd_perturbation_expansion}
    \boldsymbol{Q}(\boldsymbol{x},t) = \boldsymbol{Q}_0(\boldsymbol{x},t) + \epsilon\,\boldsymbol{Q}'(\boldsymbol{x},t) + O(\epsilon^2),
    \qquad 0 < \epsilon \ll 1,
\end{equation}
where \(\boldsymbol{Q}_0\) is a solution of the nonlinear equations, hereafter referred to as the base state, and \(\boldsymbol{Q}'\) is the perturbation about this base state. Substituting this expansion into the compact form and collecting terms of equal order in \(\epsilon\), we obtain
\begin{subequations}
    \label{eq:vd_compact_lin_continuous}
    \begin{align}
        \label{eq:vd_compact_lin_continuous_0th}
        \mathscr{M}\,\diff{\boldsymbol{Q}_0}{t} &= \mathcal{N}\!\left( \boldsymbol{Q}_0 \right), \\
        \label{eq:vd_compact_lin_continuous_1th}
        \mathscr{M}\,\diff{\boldsymbol{Q}'}{t} &= \mathcal{L}(\boldsymbol{Q}_0)\,\boldsymbol{Q}',
    \end{align}
\end{subequations}
where \(\mathcal{L}\) is the Jacobian of the nonlinear variable-density Navier--Stokes operator.

The base-state evolution and the linearized system were fully derived and characterized by \citet{kola-israel-towne-2026}; we report here only the results needed for the present analysis. The base state is a one-dimensional flow \(\boldsymbol{U}_0 = (0, 0, w_0)\) with density \(\rho_0\),
\begin{subequations}\label{eq:vd_base_state}
\begin{align}
    \rho_{0}(z,t) &= 1 + \Atw\,\erf\!\left(\frac{z}{\delta}\right), \label{eq:vd_base_state_rho0} \\[4pt]
    w_0(z,t) &= \frac{2\,\Atw}{\sqrt{\pi}\,\delta}\,
               \frac{\exp\!\left(-z^2/\delta^2\right)}{\rho_0(z,t)},
               \label{eq:vd_base_state_w0}
\end{align}
\end{subequations}
with diffusive layer thickness \(\delta(t)=\sqrt{4t/\Sch}\). Linearizing about this base state and reducing to the two-component perturbation state \(\boldsymbol{q} = (w, \rho)^\top\) gives
\begin{equation}
    \label{eq:vd_lin_continuous}
    \begin{pmatrix}
        \mathscr{B}_{ww} & \mathscr{B}_{w\rho} \\
        0 & 1
    \end{pmatrix}
    \begin{pmatrix}
        \partial_t\, w \\[4pt]
        \partial_t\,\rho
    \end{pmatrix}
    =
    \begin{pmatrix}
        \mathscr{A}_{w_0w}+\mathscr{A}_{ww} & \mathscr{A}_{w_0\rho}+\mathscr{A}_{w\rho} \\[4pt]
        \mathscr{A}_{\rho w} & \mathscr{A}_{\rho\rho}
    \end{pmatrix}
    \begin{pmatrix}
        w \\[4pt]
        \rho
    \end{pmatrix},
\end{equation}
where the operator blocks \(\mathscr{A}\) and \(\mathscr{B}\) are given in \Cref{app:operators}. The boundary conditions enforce vanishing vertical velocity and its vertical gradient at the domain boundaries, and require the density perturbation to vanish in the far field, so that the base-state stratification remains undisturbed away from the diffusive layer,
\begin{equation}
    \label{eq:vd_boundary_conditions}
    w(\pm L_z) = 0, \qquad \left.\frac{\partial w}{\partial z}\right|_{\pm L_z} = 0, \qquad
    \rho(\pm L_z) = 0.
\end{equation}
Assuming homogeneity in the horizontal directions \((x,y)\), the perturbation fields are decomposed into Fourier modes, and for the state vector \(\hat{\boldsymbol{q}}=(\hat{w},\hat{\rho})^{\top}\) the reduced system becomes
\begin{equation}
    \label{eq:vd_lti_dz_k}
    \begin{pmatrix}
        \widehat{\mathscr{B}}_{ww} & \widehat{\mathscr{B}}_{w\rho} \\
        0 & 1
    \end{pmatrix}
    \begin{pmatrix}
        \partial_t\, \hat{w} \\[4pt]
        \partial_t\,\hat{\rho}
    \end{pmatrix}
    =
    \begin{pmatrix}
        \widehat{\mathscr{A}}_{w_0w}+\widehat{\mathscr{A}}_{ww} & \widehat{\mathscr{A}}_{w_0\rho}+\widehat{\mathscr{A}}_{w\rho} \\[4pt]
        \widehat{\mathscr{A}}_{\rho w} & \widehat{\mathscr{A}}_{\rho\rho}
    \end{pmatrix}
    \begin{pmatrix}
        \hat{w} \\[4pt]
        \hat{\rho}
    \end{pmatrix},
\end{equation}
where \(\hat{w}(k,z,t)\) and \(\hat{\rho}(k,z,t)\) are the Fourier-transformed vertical velocity and density perturbations. Discretizing the inhomogeneous vertical direction with \(N_z\) points, the state at a given wavenumber is \(\hat{\mathbf{q}}_{k}(t)\in\mathbb{C}^{2N_z}\), the operators become the time-dependent matrices \(\mathsfbi{A}_{k}(t)\) and \(\mathsfbi{B}_{k}(t)\), and the linearized dynamics take the form
\begin{equation}
    \label{eq:vd_semidiscrete_system}
    \mathsfbi{B}_{k}(t)\,\difft{\hat{\mathbf{q}}_{k}}{t}
    =
    \mathsfbi{A}_{k}(t)\,\hat{\mathbf{q}}_{k},
    \qquad
    \hat{\mathbf{q}}_{k}(0)=\hat{\mathbf{q}}_{0,k},
\end{equation}
which defines the discrete propagator \(\boldsymbol{\Phi}_{t,k}\) through \(\hat{\mathbf{q}}_{k}(t)=\boldsymbol{\Phi}_{t,k}\,\hat{\mathbf{q}}_{0,k}\), with the shorthand adopted previously. Concerning the numerical solution, see \citet{kola-israel-towne-2026}, which employs spectral collocation methods with the nonlinear transformation used to cluster points within the diffusive layer.
